% GENERATED FILE. Edit canonical lessons, then run npm run generate:books.
% canonical-source-hash: fnv1a64:f772cbf97179fc18
% canonical-lessons: KA-C404-sadasya, KA-C404-dana-madu, KA-C404-protsahisu, KA-C404-samyojisu, KA-C404-sauharda, KA-R404-review-398-404

\chapter{Helping Out}
\label{ch:ka-helping-out}
\hlchaptermodality{404}

\begin{chapteropening}
\textbf{By the end of this chapter:} \emph{I can say a member, to donate, to encourage, to coordinate and goodwill.}

Complete the last lesson of chapter 404: I can say a member, to donate, to encourage, to coordinate and goodwill.
\end{chapteropening}

\section[\texorpdfstring{\kn{ಸದಸ್ಯ} (sadasya)}{ಸದಸ್ಯ (sadasya)}]{\kn{ಸದಸ್ಯ} (sadasya) — a member}

\label{lesson:KA-C404-sadasya}



\begin{quote}
Before the new one: say the Kannada for diversity, then the Kannada for to devise.
\end{quote}

\subsection*{You'll want to know: \kn{ಸದಸ್ಯ}}
\textbf{\kn{ಸದಸ್ಯ}} — \emph{sadasya} — \textquotedblleft{}a member\textquotedblright{}.

\begin{quote}
\kn{ನನ್ನ} \kn{ಅಣ್ಣ} \kn{ಈ} \kn{ಸಂಘದ} \kn{ಸದಸ್ಯ}.
\end{quote}

\begin{quote}
\emph{nanna aṇṇa ī saṅghada sadasya.}
\end{quote}

\begin{quote}
My elder brother is a member of this association.
\end{quote}

Plural: \kn{ಸದಸ್ಯರು}, \emph{sadasyaru}.

\kn{ಸದಸ್ಯತ್ವ}, membership.

\begin{cousinweb}
\textbf{\kn{ಸದಸ್ಯ}} (\emph{sadasya}) is borrowed from Sanskrit \emph{sadasya}, \textquotedblleft{}one who belongs to an assembly\textquotedblright{} (Monier-Williams), from \emph{sadas}, \textquotedblleft{}a seat, an assembly\textquotedblright{}.
\end{cousinweb}

\begin{grammarlens}[title={What you've built}]
One more word for giving reasons and opinions.
\end{grammarlens}

\subsection*{Your turn}
\begin{itemize}
\raggedright
  \item \emph{Say it:} \emph{sadasya}
  \item \emph{Say it:} \emph{sadasya}, once more
  \item \emph{Say it:} \emph{nanna aṇṇa ī saṅghada sadasya.} — My elder brother is a member of this association.
  \item \emph{Recall:} say the Kannada for diversity, then the Kannada for to devise, then say \emph{sadasya} again
\end{itemize}

\begin{tcolorbox}[breakable,colback=teal!4,colframe=teal!35!black,title={Before you move on}]
What does \kn{ಸದಸ್ಯ} mean? (\textquotedblleft{}A member\textquotedblright{}.) Say it once more.
\end{tcolorbox}
\section[\texorpdfstring{\kn{ದಾನ} \kn{ಮಾಡು} (dāna māḍu)}{ದಾನ ಮಾಡು (dāna māḍu)}]{\kn{ದಾನ} \kn{ಮಾಡು} (dāna māḍu) — to donate}

\label{lesson:KA-C404-dana-madu}



\begin{quote}
Before the new one: say the Kannada for to devise, then the Kannada for a member.
\end{quote}

\subsection*{You'll want to know: \kn{ದಾನ} \kn{ಮಾಡು}}
\textbf{\kn{ದಾನ} \kn{ಮಾಡು}} — \emph{dāna māḍu} — \textquotedblleft{}to donate\textquotedblright{}.

\begin{quote}
\kn{ನಾನು} \kn{ಆಸ್ಪತ್ರೆಗೆ} \kn{ರಕ್ತ} \kn{ದಾನ} \kn{ಮಾಡಿದೆ}.
\end{quote}

\begin{quote}
\emph{nānu āspatrege rakta dāna māḍide.}
\end{quote}

\begin{quote}
I donated blood at the hospital.
\end{quote}

Past (I …): \kn{ದಾನ} \kn{ಮಾಡಿದೆ}, \emph{dāna māḍide}.

\kn{ದಾನ} alone is a gift or a donation.

\begin{cousinweb}
\textbf{\kn{ದಾನ} \kn{ಮಾಡು}} (\emph{dāna māḍu}) is built from Sanskrit \emph{dāna}, \textquotedblleft{}giving, a gift\textquotedblright{} (CDIAL 6265), with \emph{māḍu}, \textquotedblleft{}to do\textquotedblright{}.
\end{cousinweb}

\begin{grammarlens}[title={What you've built}]
One more word for giving reasons and opinions.
\end{grammarlens}

\subsection*{Your turn}
\begin{itemize}
\raggedright
  \item \emph{Say it:} \emph{dāna māḍu}
  \item \emph{Say it:} \emph{dāna māḍu}, once more
  \item \emph{Say it:} \emph{nānu āspatrege rakta dāna māḍide.} — I donated blood at the hospital.
  \item \emph{Recall:} say the Kannada for to devise, then the Kannada for a member, then say \emph{dāna māḍu} again
\end{itemize}

\begin{tcolorbox}[breakable,colback=teal!4,colframe=teal!35!black,title={Before you move on}]
What does \kn{ದಾನ} \kn{ಮಾಡು} mean? (\textquotedblleft{}To donate\textquotedblright{}.) Say it once more.
\end{tcolorbox}
\section[\texorpdfstring{\kn{ಪ್ರೋತ್ಸಾಹಿಸು} (prōtsāhisu)}{ಪ್ರೋತ್ಸಾಹಿಸು (prōtsāhisu)}]{\kn{ಪ್ರೋತ್ಸಾಹಿಸು} (prōtsāhisu) — to encourage}

\label{lesson:KA-C404-protsahisu}



\begin{quote}
Before the new one: say the Kannada for a member, then the Kannada for to donate.
\end{quote}

\subsection*{You'll want to know: \kn{ಪ್ರೋತ್ಸಾಹಿಸು}}
\textbf{\kn{ಪ್ರೋತ್ಸಾಹಿಸು}} — \emph{prōtsāhisu} — \textquotedblleft{}to encourage\textquotedblright{}.

\begin{quote}
\kn{ಅಪ್ಪ} \kn{ನನ್ನನ್ನು} \kn{ಆಟದಲ್ಲಿ} \kn{ಪ್ರೋತ್ಸಾಹಿಸಿದರು}.
\end{quote}

\begin{quote}
\emph{appa nannannu āṭadalli prōtsāhisidaru.}
\end{quote}

\begin{quote}
Father encouraged me in sports.
\end{quote}

Past (I …): \kn{ಪ್ರೋತ್ಸಾಹಿಸಿದೆ}, \emph{prōtsāhiside}.

\kn{ಪ್ರೋತ್ಸಾಹ}, encouragement.

\begin{cousinweb}
\textbf{\kn{ಪ್ರೋತ್ಸಾಹಿಸು}} (\emph{prōtsāhisu}) is built on Sanskrit \emph{protsāha}, \textquotedblleft{}great zeal, stimulus\textquotedblright{} (Monier-Williams), with the Kannada verb ending -\emph{isu}.
\end{cousinweb}

\begin{grammarlens}[title={What you've built}]
One more word for giving reasons and opinions.
\end{grammarlens}

\subsection*{Your turn}
\begin{itemize}
\raggedright
  \item \emph{Say it:} \emph{prōtsāhisu}
  \item \emph{Say it:} \emph{prōtsāhisu}, once more
  \item \emph{Say it:} \emph{appa nannannu āṭadalli prōtsāhisidaru.} — Father encouraged me in sports.
  \item \emph{Recall:} say the Kannada for a member, then the Kannada for to donate, then say \emph{prōtsāhisu} again
\end{itemize}

\begin{tcolorbox}[breakable,colback=teal!4,colframe=teal!35!black,title={Before you move on}]
What does \kn{ಪ್ರೋತ್ಸಾಹಿಸು} mean? (\textquotedblleft{}To encourage\textquotedblright{}.) Say it once more.
\end{tcolorbox}
\section[\texorpdfstring{\kn{ಸಂಯೋಜಿಸು} (saṃyōjisu)}{ಸಂಯೋಜಿಸು (saṃyōjisu)}]{\kn{ಸಂಯೋಜಿಸು} (saṃyōjisu) — to coordinate; to combine}

\label{lesson:KA-C404-samyojisu}



\begin{quote}
Before the new one: say the Kannada for to donate, then the Kannada for to encourage.
\end{quote}

\subsection*{You'll want to know: \kn{ಸಂಯೋಜಿಸು}}
\textbf{\kn{ಸಂಯೋಜಿಸು}} — \emph{saṃyōjisu} — \textquotedblleft{}to coordinate; to combine\textquotedblright{}.

\begin{quote}
\kn{ಅವಳು} \kn{ಶಾಲೆಯ} \kn{ಕಾರ್ಯಕ್ರಮವನ್ನು} \kn{ಸಂಯೋಜಿಸಿದಳು}.
\end{quote}

\begin{quote}
\emph{avaḷu śāleya kāryakramavannu saṃyōjisidaḷu.}
\end{quote}

\begin{quote}
She coordinated the school event.
\end{quote}

Past (I …): \kn{ಸಂಯೋಜಿಸಿದೆ}, \emph{saṃyōjiside}.

\kn{ಸಂಯೋಜಕ}, a coordinator.

\begin{cousinweb}
\textbf{\kn{ಸಂಯೋಜಿಸು}} (\emph{saṃyōjisu}) is built on Sanskrit \emph{saṃyojana}, \textquotedblleft{}the act of joining or uniting\textquotedblright{} (Monier-Williams), with the Kannada verb ending -\emph{isu}.
\end{cousinweb}

\begin{grammarlens}[title={What you've built}]
One more word for giving reasons and opinions.
\end{grammarlens}

\subsection*{Your turn}
\begin{itemize}
\raggedright
  \item \emph{Say it:} \emph{saṃyōjisu}
  \item \emph{Say it:} \emph{saṃyōjisu}, once more
  \item \emph{Say it:} \emph{avaḷu śāleya kāryakramavannu saṃyōjisidaḷu.} — She coordinated the school event.
  \item \emph{Recall:} say the Kannada for to donate, then the Kannada for to encourage, then say \emph{saṃyōjisu} again
\end{itemize}

\begin{tcolorbox}[breakable,colback=teal!4,colframe=teal!35!black,title={Before you move on}]
What does \kn{ಸಂಯೋಜಿಸು} mean? (\textquotedblleft{}To coordinate; to combine\textquotedblright{}.) Say it once more.
\end{tcolorbox}
\section[\texorpdfstring{\kn{ಸೌಹಾರ್ದ} (sauhārda)}{ಸೌಹಾರ್ದ (sauhārda)}]{\kn{ಸೌಹಾರ್ದ} (sauhārda) — goodwill; harmony}

\label{lesson:KA-C404-sauharda}



\begin{quote}
Before the new one: say the Kannada for to encourage, then the Kannada for to coordinate.
\end{quote}

\subsection*{You'll want to know: \kn{ಸೌಹಾರ್ದ}}
\textbf{\kn{ಸೌಹಾರ್ದ}} — \emph{sauhārda} — \textquotedblleft{}goodwill; harmony\textquotedblright{}.

\begin{quote}
\kn{ಎರಡು} \kn{ಊರುಗಳ} \kn{ನಡುವೆ} \kn{ಸೌಹಾರ್ದ} \kn{ಇದೆ}.
\end{quote}

\begin{quote}
\emph{eraḍu ūrugaḷa naḍuve sauhārda ide.}
\end{quote}

\begin{quote}
There is goodwill between the two towns.
\end{quote}

\kn{ಸೌಹಾರ್ದದಿಂದ}, amicably.

\begin{cousinweb}
\textbf{\kn{ಸೌಹಾರ್ದ}} (\emph{sauhārda}) is borrowed from Sanskrit \emph{sauhārda}, \textquotedblleft{}good-heartedness, friendship\textquotedblright{} (Monier-Williams), from \emph{suhṛd}, \textquotedblleft{}a friend\textquotedblright{}.
\end{cousinweb}

\begin{grammarlens}[title={What you've built}]
One more word for giving reasons and opinions.
\end{grammarlens}

\subsection*{Your turn}
\begin{itemize}
\raggedright
  \item \emph{Say it:} \emph{sauhārda}
  \item \emph{Say it:} \emph{sauhārda}, once more
  \item \emph{Say it:} \emph{eraḍu ūrugaḷa naḍuve sauhārda ide.} — There is goodwill between the two towns.
  \item \emph{Recall:} say the Kannada for to encourage, then the Kannada for to coordinate, then say \emph{sauhārda} again
\end{itemize}

\begin{tcolorbox}[breakable,colback=teal!4,colframe=teal!35!black,title={Before you move on}]
What does \kn{ಸೌಹಾರ್ದ} mean? (\textquotedblleft{}Goodwill; harmony\textquotedblright{}.) Say it once more.
\end{tcolorbox}
\section[(dialogue)]{Looking back: from Protest and Struggle to Helping Out}

\label{lesson:KA-R404-review-398-404}



\begin{quote}
Say the Kannada for to protest and a struggle.
\end{quote}

\subsection*{Your turn}
Say these aloud:

\begin{itemize}
\raggedright
  \item to protest, a struggle, a petition, to mobilise, an activist
  \item poverty, a labourer, a privilege, to distribute, rural
  \item a budget, income, a subsidy, to utilise, to waste
  \item development, economic, finance, to promote, to invest
  \item corruption, a bribe, transparent, to conceal, to recommend
  \item society, self-reliance, awareness, diversity, to devise
  \item a member, to donate, to encourage, to coordinate, goodwill
\end{itemize}

\begin{tcolorbox}[breakable,colback=teal!4,colframe=teal!35!black,title={Before you move on}]
To protest? (\textbf{\kn{ಪ್ರತಿಭಟಿಸು}}, \emph{pratibhaṭisu}.) A struggle? (\textbf{\kn{ಹೋರಾಟ}}, \emph{hōrāṭa}.) A petition? (\textbf{\kn{ಮನವಿ}}, \emph{manavi}.) To mobilise? (\textbf{\kn{ಸಂಘಟಿಸು}}, \emph{saṅghaṭisu}.) An activist? (\textbf{\kn{ಕಾರ್ಯಕರ್ತ}}, \emph{kāryakarta}.) Poverty? (\textbf{\kn{ಬಡತನ}}, \emph{baḍatana}.) A labourer? (\textbf{\kn{ಕಾರ್ಮಿಕ}}, \emph{kārmika}.) A privilege? (\textbf{\kn{ಸವಲತ್ತು}}, \emph{savalattu}.) To distribute? (\textbf{\kn{ವಿತರಿಸು}}, \emph{vitarisu}.) Rural? (\textbf{\kn{ಗ್ರಾಮೀಣ}}, \emph{grāmīṇa}.) A budget? (\textbf{\kn{ಬಜೆಟ್}}, \emph{bajeṭ}.) Income? (\textbf{\kn{ಆದಾಯ}}, \emph{ādāya}.) A subsidy? (\textbf{\kn{ಸಬ್ಸಿಡಿ}}, \emph{sabsiḍi}.) To utilise? (\textbf{\kn{ವಿನಿಯೋಗಿಸು}}, \emph{viniyōgisu}.) To waste? (\textbf{\kn{ವ್ಯರ್ಥ} \kn{ಮಾಡು}}, \emph{vyartha māḍu}.) Development? (\textbf{\kn{ಅಭಿವೃದ್ಧಿ}}, \emph{abhivṛddhi}.) Economic? (\textbf{\kn{ಆರ್ಥಿಕ}}, \emph{ārthika}.) Finance? (\textbf{\kn{ಹಣಕಾಸು}}, \emph{haṇakāsu}.) To promote? (\textbf{\kn{ಉತ್ತೇಜಿಸು}}, \emph{uttējisu}.) To invest? (\textbf{\kn{ಹೂಡು}}, \emph{hūḍu}.) Corruption? (\textbf{\kn{ಭ್ರಷ್ಟಾಚಾರ}}, \emph{bhraṣṭācāra}.) A bribe? (\textbf{\kn{ಲಂಚ}}, \emph{lañca}.) Transparent? (\textbf{\kn{ಪಾರದರ್ಶಕ}}, \emph{pāradarśaka}.) To conceal? (\textbf{\kn{ಮರೆಮಾಚು}}, \emph{maremācu}.) To recommend? (\textbf{\kn{ಶಿಫಾರಸು} \kn{ಮಾಡು}}, \emph{śiphārasu māḍu}.) Society? (\textbf{\kn{ಸಮಾಜ}}, \emph{samāja}.) Self-reliance? (\textbf{\kn{ಸ್ವಾವಲಂಬನೆ}}, \emph{svāvalambane}.) Awareness? (\textbf{\kn{ಜಾಗೃತಿ}}, \emph{jāgṛti}.) Diversity? (\textbf{\kn{ವೈವಿಧ್ಯ}}, \emph{vaividhya}.) To devise? (\textbf{\kn{ರೂಪಿಸು}}, \emph{rūpisu}.) A member? (\textbf{\kn{ಸದಸ್ಯ}}, \emph{sadasya}.) To donate? (\textbf{\kn{ದಾನ} \kn{ಮಾಡು}}, \emph{dāna māḍu}.) To encourage? (\textbf{\kn{ಪ್ರೋತ್ಸಾಹಿಸು}}, \emph{prōtsāhisu}.) To coordinate? (\textbf{\kn{ಸಂಯೋಜಿಸು}}, \emph{saṃyōjisu}.) Goodwill? (\textbf{\kn{ಸೌಹಾರ್ದ}}, \emph{sauhārda}.)
\end{tcolorbox}
